\section{Upper-bound classification under the recognition hypothesis}\label{sec:upper-bound-classification}



The lower bounds of Section~\ref{sec:lower-bounds} rule out the
trivialization in which some primitive role lacks an admissible
witness in \FATCD{}. The upper-bound direction addresses the
complementary question: whether some additional, seventh closure role
is needed in the covered scope. This section shows that, under the
recognition hypothesis of Definition~\ref{def:recognition-hypothesis}, every
kind residual in the covered scope, every bookkeeping annotation, and every record residual of every
well-formed decomposition record, is licensed only labels among (i)--(ix) of
Definition~\ref{def:residual-scheme}, and that the recognized forms of
probability, causality and agency receive such labels without it. The argument proceeds in three steps. First, we state the
alignment rules that constrain how raw features may project to
\Pone{}, \ldots, \Psix{}. Second, we close the three high-priority
candidate seventh-role threats: probability and uncertainty, residual
causality, and residual agency. Third, we record the upper-bound
classification theorem that collects these local classifications into
a single covered-scope statement.

\subsection{Role-projection alignment rules}\label{subsec:alignment-rules}

The role meanings of Definition~\ref{def:primitive-roles} carry
alignment rules constraining which raw features may be interpreted as
active projections of each role. These rules refine the
alignments of Remark~\ref{rem:repaired-alignments} and form the
discipline under which the classification theorem is proved.

\begin{proposition}[Role-projection alignment]\label{prop:alignment-rules}
Under honest bookkeeping and the witness discipline of
Definition~\ref{def:role-projection}, each role projection is
constrained as follows.
\begin{itemize}[topsep=2pt,itemsep=1pt]
  \item \Pone{}: admissible only when the candidate cluster carries a map or partial-map record of declared use update, source and target relation records of declared use quotient, and a certificate of alternative (a) or (b) of the \(W_1\) row. In particular, a cluster with no map or partial-map record
    of declared use update, or no relation record of declared use quotient,
    does not satisfy \(W_1\).
  \item \Ptwo{}: admissible only when the candidate carries a macro
    specification, a realization relation, specification-satisfaction
    data, and evidence of certified unrealizability or failure of
    representability over the declared packaging kernel; when a
    satisfying realizer exists, the candidate must also carry a comparison
    with the classes of the kernel of a declared packaging map, whose
    declared domain contains the declared candidate class, certifying
    that some class of that kernel meets both the realizer set and its
    complement in the candidate class.
    Macro-level predicates or objectives count only insofar as they are
    treated as content to be represented. When no realizer exists, the candidate class must be the full declared domain of a map record of the cluster of declared use packaging or update, and the cluster must record a true evaluation of the specification at a declared same-kind record outside that class. Mere quotient equality and
    arbitrary relations without representability structure are not
    admissible as \Ptwo{} projections.
  \item \Pthree{}: admissible only when the candidate carries a route
    grammar (generating moves and all rewrite rules recorded in \(D\)), generated same-source, same-target route terms, and
    route-comparison data with a non-reducibility certificate of either
alternative in the \(W_3\) row. Generic transitions, generic causal arrows, generic
    path records, and route undefinedness arising from a partial update
    map are not admissible as \Pthree{} projections unless they are
    two defined same-source, same-target route records carrying a
    non-reducibility certificate of either alternative in the \(W_3\) row.
  \item \Pfour{}: admissible only when the candidate carries a stage,
    order, transition or refinement record (Definition~\ref{def:raw-record})
    with a law in the sense of the witness table (for a transition record,
    a certificate that its table is defined on its declared domain with
    values in its declared codomain suffices). Bare labels, arbitrary
    relations without order, transition, or refinement law, quotient
    packaging, and audit bookkeeping are not admissible as \Pfour{}
    projections.
  \item \Pfive{}: admissible only when the candidate carries a quotient
    relation, a packaging map, the identification of distinct
    micro-elements, an object record declared as the quotient object, and
    quotient-validity evidence. Generic structural paths, derivations,
    topologies or geometries merely because they have paths, and causal
    chains are not admissible as \Pfive{} projections unless quotient
    or packaging data is present.
  \item \Psix{}: admissible only when the cluster carries a declared event,
    a trace record, a record tied to the event that is that trace record or a
    provenance record also tied to that trace record, a replay or check record tied to the trace carrying a certificate of a replay or provenance property of
    the trace, and \(\mathrm{Audit}_0(D)\) contains an entry \(a\), not necessarily a cluster member, referring to a claim
    annotation, to these event and trace records and to that replay or check
    record, whose conclusion is the formula of that certificate.
    Passive logs alone, arbitrary metadata, and certificate-only
    content without audit/provenance/replay/check structure are not
    admissible as \Psix{} projections.
\end{itemize}
\end{proposition}

\begin{proof}
By Definition~\ref{def:role-projection}, an admissible projection
\(\pi_i(D,C)\) requires \(W_i(C)\), and \(W_i\) is defined by the row for
\(P_i\) of the \emph{Witness types} table in
Section~\ref{sec:primitive-role-architecture}. Each constraint below follows
by reading off the record kinds that row requires; the table is the
definition and this proposition states its consequences for candidate
features. We go through the roles in turn.

For \Pone{}, the fixed meaning is update-vs-quotient compatibility and
descent obstruction. The witness type for that meaning is exactly a
selected update together with source and target relation records of declared use quotient and a
certificate of descent compatibility, descent obstruction, or
update-vs-quotient mismatch; absent such a witness the candidate
exhibits no descent datum to compare, so it cannot be an active \Pone{}
projection. In particular generic semantics, generic predicate meaning,
generic selection criteria, and generic agency carry no descent witness
and are inadmissible, which is the alignment of
Remark~\ref{rem:repaired-alignments}.

For \Ptwo{}, the fixed meaning is macro-specification representability.
Its witness type is a macro specification, a realization relation,
specification-satisfaction data, and a certificate of unrealizability or of
failure of representability over the declared packaging kernel; a macro-level predicate or objective
enters only when treated as content to be represented, since only then
does it instantiate the representability relation. Mere quotient
equality is \Pfive{} data, an arbitrary relation lacking
representability structure supplies no realization relation, and
empirical observations lacking level and threshold data fail the witness
discipline; none is an admissible \Ptwo{} projection.

For \Pthree{}, the fixed meaning is enriched route/coherence mismatch,
and \Pthree{} is not a directionality certificate. Its witness type is
a route grammar (generating moves and all rewrite rules recorded in \(D\)), generated same-source same-target route terms,
and route-comparison data carrying a non-reducibility certificate of either
alternative in the \(W_3\) row. A candidate without such a route
comparison has nothing of the \Pthree{} type to compare. Generic
transitions and generic causal arrows are \Pfour{} data, generic path
records are not route-comparison data, and route undefinedness arising
from a partial update map is \Pone{} content (Table~\ref{tab:boundary-matrix},
row \Pone{}/\Pthree{}); none projects to
\Pthree{} unless the route/coherence threshold is explicitly present and
the comparison is specifically of the \Pthree{} type. This is the
alignment that generic transition is not \Pthree{}.

For \Pfour{}, the fixed meaning is stage, order, transition, and
refinement structure. Its witness type is a stage, order, transition or
refinement record carrying a law beyond bare labels; staged dynamics and
regime ordering qualify when recorded in one of these kinds. Bare labels, arbitrary relations without an order,
transition, or refinement law, quotient packaging (which is \Pfive{}
data), and audit bookkeeping (which is \Psix{} data) supply no such law
and are inadmissible as \Pfour{} projections.

For \Pfive{}, the fixed meaning is quotient packaging. Its witness type
is a quotient relation, a packaging map, the identification of distinct
micro-elements, an object record declared as the quotient object, and
quotient-validity evidence. Generic structural paths, derivations,
topologies or geometries that merely carry paths, and causal chains
supply no quotient relation or packaging map, so they are inadmissible
absent quotient or packaging data; this is the alignment that
generic path or derivation is not \Pfive{}.

For \Psix{}, the fixed meaning is audit, provenance, replay, and
checkability. Its witness type is a declared event, a trace record tied to it directly or through a provenance record, a replay or check record carrying a certified replay or provenance property of the trace, and an entry of \(\mathrm{Audit}_0(D)\) concluding that property. Passive logs without check
structure, arbitrary metadata, and certificate-only content lacking
audit/provenance/replay/check structure supply no checkable audit datum
and are inadmissible as \Psix{} projections.

In each case the admissibility constraint is forced by the fixed role
meaning through the witness discipline: a projection is admissible only
with the witness type that the meaning of the role fixes. Formally, each
bullet is the contrapositive of the row conjunct it names: \(W_1\) requires an
update map and quotient relations; \(W_2\) a macro-specification and a
realization relation; \(W_3\) two defined parallel routes and a certificate;
\(W_4\) a law-bearing record; \(W_5\) a quotient relation, a packaging map and a
quotient object; and \(W_6\) an event, a trace, a certified check and an
\(\mathrm{Audit}_0\) entry. The remaining sentences gloss role meanings and are
not used.
\end{proof}

\begin{remark}[Three excluded alignments, held fixed]\label{rem:three-repaired-alignments}
Proposition~\ref{prop:alignment-rules} fixes three alignments that
are held throughout: generic transition data does not project to
\Pthree{}, generic path or derivation data does not project to
\Pfive{}, and generic semantic or interpretation data does not project
to \Pone{}. These are carried into every step of the classification
theorem below.
\end{remark}

\subsection{Probability, causality and agency}\label{subsec:threat-closures}

Three classes of candidate seventh-role feature require particular
scrutiny: probability and uncertainty, residual causality in the
sense of the interventional and counterfactual
hierarchy~\citep{Pearl2009}, and residual agency in the planning
sense~\citep{Bratman1987}. Each is classified in a proposition. In the titles of
Propositions~\ref{prop:probability-closure}--\ref{prop:agency-closure},
\emph{closure} of a threat means only that each listed form is licensed a label
among (i)--(ix); it is unrelated to closure operations, to the condition
\(\Closed{D}\), or to the closure content of
Definition~\ref{def:recognition-hypothesis}.
In every case the candidate is classified under the alignment rules
of Proposition~\ref{prop:alignment-rules} together with the residual
scheme of Definition~\ref{def:residual-scheme}; no candidate is
dismissed by definition. These three come first because they are the
most natural candidates for a seventh role; the general classification
theorem follows. A \emph{probability or uncertainty feature} is a residual
feature with a field of declared probability, weight, kernel, confidence or
error type, or an empirical-bridge annotation whose identified bridge claim has
such a field in its \(\Audit{D}\) entry; a \emph{residual-causality feature} is one with a field of declared
intervention, counterfactual or causal-dependence type; and a
\emph{residual-agency feature} is one with a field of declared goal,
objective, plan, policy or choice type, or an empirical-bridge annotation
whose identified bridge claim has such a field in its \(\Audit{D}\) entry.
Value-type tags enter no label criterion (Definition~\ref{def:raw-record}), and
the proofs below case only on content; hence each proposition below holds for
every kind residual whose content takes one of its forms, whether or not it
carries such a tag. Each recognized form is specified so that a criterion
among (i)--(ix) applies to it; the propositions list which presentations of
probabilistic, causal and agentive content the witness rows, seed clauses
and annotation kinds already accommodate.

\paragraph{The \(P_i\)-form.}\label{par:pi-form} In these propositions, a feature \(r\) has \emph{\(P_i\)-form} when \(r\) is a
record of a kind occurring in the \(W_i\) row of the witness table,
contributes a required field to a finite cluster \(C\) satisfying \(W_i(C)\),
and every predicate, comparison, invariant, law or claim field of \(r\) is
required for \(W_i\) at \(C\) or is an input to its check
(Section~\ref{subsec:residual-classification}); or when \(r\) carries a covered
\(P_i\) seed and each of its predicate, comparison, invariant, law or claim fields is a named seed field of that clause. Such
a feature satisfies criterion (i) when that criterion holds, and criterion (v)
otherwise. In the bullets below, \(P_i\)-aligned content is
labelled (i) or (v); a bullet does not assert that the decomposition record
names an active projection for it. A seed alone is not an active projection.
\par\emph{Examples.} For instance, a transition record whose fields are its endpoints, a kernel
declared with values in the distributions on its state set, and the law that
each row is such a distribution has \(P_4\)-form: in an aligned cluster in which it is the only law-bearing
record, deleting the row law leaves the cluster without a law, so \(W_4\)
fails. By contrast, the predicate \(N\) in the original \(D_{\mathrm{norm}}\) of
Remark~\ref{rem:recognition-open} has no \(P_i\)-form; neither does a
standalone transition record without a certified law or covered seed. For a
composite form, the feature is carried by a finite set \(R_f\) of records, and
each record of \(R_f\) has the \(P_i\)-form of its own component role (so, by
the second sentence of this paragraph, it satisfies criterion (i) or (v), and
also (ii) when it contributes required fields to two named active clusters and every
predicate, comparison, invariant, law or claim field is required for, or an
input to the check of, either cluster's witness predicate).
\par\emph{Composite forms.} For a composite form, ``classifies as a composite'' is a derived
reading of the finite record set \(R_f\), not an additional residual
label: \(\operatorname{Classify}\) still labels each member of \(R_f\).
Label (ii) applies to a member only under the two-cluster condition just
stated.
\par\emph{How the propositions proceed.} Each proposition takes a kind residual of \(D\) whose content
is of the named kind --- probability or uncertainty; causation in the
interventional and counterfactual sense; agency in the planning sense --- and
passes it through the residual scheme under the alignment rules. Its proof
checks, for each recognized form, that the feature supplies the witness type
of the role it is assigned to, or is an annotation or an outside-domain
structure; the causality and agency propositions also show that one role is
unavailable: \Pfive{} without quotient packaging, and \Pone{} without an
update-vs-quotient descent witness. A stochastic transition record carrying a certified row law, with no assertion
fields beyond the named fields of its covered \(P_4\) seed, has \(P_4\)-form
whether or not attachment data are present. Relative to a fixed well-formed
\(\Dec{D}\), it satisfies criterion (i) if that criterion holds, and criterion
(v) otherwise; an uncertainty record with provenance and replay data
has \(P_6\)-form when each of its predicate, comparison, invariant, law or claim
fields is required for, or an input to, a \(W_6\) check, or is a named field of
its seed; an empirical-bridge annotation attaching a measurement error to a substrate observable is licensed label (vii) when it meets criterion (vii), the error record itself being classified by its own fields; and a
confidence level recorded as the threshold predicate of a threshold annotation
with no other field besides its typing, target and label fields, all of whose
read records satisfy (R1), is an activation-threshold refinement. A normalization predicate with no free variable, such as the
one in \(D_{\mathrm{norm}}\) of Remark~\ref{rem:recognition-open}, is not a
\Ptwo{} seed merely because it states that finitely many weights sum to
one. That remark shows that this
specified predicate receives label (x), and discusses conservation laws and
emergent obstructions separately.
A record of \(P_i\)-form satisfies alternative (R1) of
Definition~\ref{def:recognition-hypothesis}, the threshold and bridge forms below
satisfy (R2), and the outside-domain forms satisfy (R5); the propositions below
are thus the cases, for these three classes, of the pointwise clause of
Theorem~\ref{thm:upper-bound-classification}, and assume nothing about the other
features of \(D\).

\begin{proposition}[Closure of probability and uncertainty]\label{prop:probability-closure}
\textup{(a)} Every kind residual (Definition~\ref{def:residual-scheme}), in particular every probability or uncertainty feature, of a description
\(D \in \FATCD{}\) whose content takes one of the recognized forms below,
where a form naming a \(P_i\) projection or unabsorbed \(P_i\)-aligned content
requires the feature to have \(P_i\)-form (paragraph \emph{The \(P_i\)-form}
before Proposition~\ref{prop:probability-closure}),
is licensed, relative to every well-formed decomposition record \(\Dec{D}\), at least one of the labels of
\begin{itemize}[topsep=2pt,itemsep=1pt]
  \item \(P_2\)-aligned content, label (i) or (v), when the feature has \(P_2\)-form (for
    instance a macro-specification with one free variable and a realization
    relation whose recorded evaluations take both truth values);
  \item \(P_4\)-aligned content, label (i) or (v), when the feature has \(P_4\)-form (for
    instance a stochastic transition record carrying its row law);
  \item \(P_6\)-aligned content, label (i) or (v), when the feature has \(P_6\)-form (for
    instance an audited uncertainty-bookkeeping record);
  \item an activation-threshold refinement, when the feature is a threshold
    annotation targeting a cluster named for an active projection, with a
    confidence level recorded as its threshold predicate
    (Definition~\ref{def:threshold-attachment}) as its only field besides
    typing, target and label, and that predicate reads at least one record, every record whose values it reads is listed in the annotation's target field, and each such record satisfies (R1); relative to a record that names its target cluster for
    no active projection, the annotation satisfies criterion (ix-b);
  \item an empirical-bridge annotation, when the feature is itself an empirical-bridge annotation identifying an admissible bridge claim (Definitions~\ref{prop:honest-bookkeeping} and~\ref{prop:empirical-bridge}) whose \(\Audit{D}\) entry carries the five items of Definition~\ref{prop:empirical-bridge}, with no fields beyond its typing, targets and that identification;
  \item an outside-domain structure, label (viii), when an assertion field of the feature uses a declared predicate whose extension \(D\) does not record or quantifies over a sort whose carrier \(D\) does not record, the asserted property is true in some expansion of \(D\) and false in another, and every other assertion field of that record either also meets the outside-scope condition or is required for \(W_i\) at a cluster aligned with \(P_i\), an input to its check at such a cluster, or named in a seed clause that the record carries (for instance the \(\mathrm{IsDist}\) presentation of Remark~\ref{rem:recognition-open});
  \item label (i), (ii) or (v), as a member of a composite of the foregoing
    projections, when the feature is one of a finite set of records each of which has the \(P_i\)-form of one of the roles above (so, by the paragraph
    \emph{The \(P_i\)-form}, it satisfies criterion (i) or (v), and also (ii)
    when it contributes required fields to two named active clusters and every
    predicate, comparison, invariant, law or claim field is required for, or an
    input to the check of, either cluster's witness predicate);
  \item unabsorbed role-aligned content, label~(v) of
    Definition~\ref{def:residual-scheme}, when it has \(P_i\)-form but
    criterion (i) does not hold for it, for instance because the cluster lacks
    the attachment data of Definition~\ref{def:threshold-attachment}.
\end{itemize}
A content record targeted by such an empirical-bridge annotation is not thereby classified; it is classified by its own fields and is licensed label (ix-c) exactly when either branch of that criterion holds.
Each case licenses a label among (i)--(ix) of
Definition~\ref{def:residual-scheme}, so neither (x) nor (xi) is licensed for
the feature. Probability or uncertainty content of any of these recognized forms therefore
requires no seventh role; content of none of these forms is governed by the
recognition hypothesis (Definition~\ref{def:recognition-hypothesis}).
\textup{(b)} Let a content record of \(D\) have as its only assertion field the formula \(\bigwedge_k v(w_k)\ge0\wedge\sum_k v(w_k)=1\) over value fields of pairwise distinct declared records \(w_1,\dots,w_n\), none of which is the replaced record, suppose the recorded values \(v(w_k)\) are nonnegative and sum to one (so the recomputed row check is true), and let \(D'\) be the description obtained from \(D\) by replacing that
feature, with every incoming reference and affected annotation retargeted in a
type-correct way and with the resulting description \(D'\) in \FATCD{} (in particular, every realization relation still lists exactly the members of its candidate class at which its target's formula is true in \(D'\)), by a transition record \(K:\{*\}\to\mathrm{Dist}(\{w_1,\dots,w_n\})\) (here \(K\) is a transition record, not a candidate class)
whose declared source is an object record \(\{*\}\) listing a declared state
record \(*\), whose declared target is \(\mathrm{Dist}(X)\) for an object record
\(X\) listing \(w_1,\dots,w_n\), recorded by a codomain reference to \(X\) and a
datum field marking the codomain as distribution-valued (paragraph
\emph{Laws}), whose unique row is the list of terms
\(v(w_1),\dots,v(w_n)\), and whose law certificate is
\((\bigwedge_k v(w_k)\ge0\wedge\sum_k v(w_k)=1,\mathrm{true})\); the records
\(*\), \(\{*\}\) and \(X\) are added to \(D'\).
Then \(K\) is a covered \Pfour{} seed of \(D'\) and has \(P_4\)-form, so
relative to every well-formed decomposition record of \(D'\) it is licensed
label (i) or (v), and neither (x) nor (xi) (its only assertion field is the
law certificate, a named seed field of clause (4), so \(K\) has \(P_4\)-form by
the second branch of the paragraph \emph{The \(P_i\)-form}); the added records \(*\),
\(\{*\}\) and \(X\) have no assertion field and satisfy (R1) or (R3). The \(w_i\) and their values are
unchanged, source and target declarations are added, and the claim and audit
annotations of the replaced feature are retargeted to the row-law check of
\(K\); in \(D\) itself the feature is not thereby classified.
\end{proposition}

\begin{proof}
Let \(r\) be a probability or uncertainty feature (a kind residual) of
\(D \in \FATCD{}\), and let it pass through the residual scheme of
Definition~\ref{def:residual-scheme} under the alignment rules of
Proposition~\ref{prop:alignment-rules}. Role-aligned classifications below require the \(P_i\)-form specified in the
proposition. Such a record satisfies criterion (i) when that criterion holds and
criterion (v) otherwise. Composite forms are determined record by record;
criterion (ii) applies under the two-cluster condition of the paragraph
\emph{The \(P_i\)-form}. We case on the content of \(r\).

If the probabilistic content is a macro specification or a
representability or non-representability condition on a description that
contributes a required field to a cluster satisfying \(W_2\), with its
realization relation and certificate, then it supplies the \Ptwo{}
witness, so it has \(P_2\)-form and receives label (i) or (v) (paragraph
\emph{The \(P_i\)-form}). If instead the content is a stochastic transition,
regime, or staged structure carrying a valid stage, order, transition,
or refinement law and has \(P_4\)-form, then it receives label (i) or (v); the
row law of a kernel is such a law, as in the example before this proposition,
and a normalization assertion that is not required for \(W_4\) is carried by a
separate predicate record, classified in its own right
(Remark~\ref{rem:recognition-open}). If the content is an audited record of uncertainty
bookkeeping --- event reference, provenance, replay or check structure
--- it has \(P_6\)-form. A \(P_6\)-form record receives (i) when its named
active cluster carries the certified trace check and honest audit required by
\(W_6\); a covered seed without that full witness receives (v). If the
feature is a threshold annotation of a cluster named for an active projection
whose only field besides typing, target and label is a threshold predicate recording a confidence level
that decides when the projection is attempted, and that predicate reads at least one record, every record whose values it reads is listed in the annotation's target field, and each such record satisfies (R1), then it adds no role content and receives
label (v-b): it is an activation-threshold
refinement. Relative to a record that names its target cluster for no active projection,
the annotation is a structural annotation whose fields are its typing, target
and threshold predicate, with no independent role-theoretic assertion, and
(v), (vi) and (vii) do not apply, so it satisfies criterion (ix-b). If it is an empirical-bridge annotation whose bridge claim is admissible and has an audit
entry with the five items of Definition~\ref{prop:empirical-bridge}, and with
no fields beyond its typing, targets and that identification, it
receives (vii). A content record targeted by such an annotation is not thereby classified; it is classified by its own fields and is licensed label (ix-c) exactly when either branch of that criterion holds. Finally, if several role aspects occur together, take their finite
record set \(R_f\). By the composite-form hypothesis each member has
its component role's \(P_i\)-form, so it satisfies criterion (i) or (v), and
also (ii) under the two-cluster condition stated there;
the set has the composite reading of those classifications. These are the recognized forms of probability and uncertainty content, and
each reduces to the six roles as stated. A probability feature whose content
has none of these forms is not classified by this proposition. In particular,
the bare normalization predicate of Remark~\ref{rem:recognition-open} receives
label (x); other such constraints must be checked against the recognition
alternatives of Definition~\ref{def:recognition-hypothesis}.
A feature whose probabilistic content lies outside the covered scope is classified as an outside-domain structure by Definition~\ref{def:residual-scheme}, not absorbed by default: in the outside-domain case, the asserted property is true in some expansion of \(D\) and false in another, and together with the stated condition on the remaining assertion fields this satisfies criterion (viii). The
closure is therefore scoped, as stated.

For (b), the row of \(K\) is a map from the declared one-point
domain to the declared set \(\mathrm{Dist}(\{w_1,\dots,w_n\})\), and the finite
check of nonnegativity and sum one certifies that its value lies in that
codomain. This is a law in the sense of the witness table, so \(K\) is a
law-bearing transition record, which is a covered \Pfour{} seed by
Definition~\ref{def:channel-status}. Its declared source and target records are fixed by reference fields, and the
value fields of the \(w_k\) are read by that check; none of these records carries
an assertion of its own.
Reference closure and honest bookkeeping of \(D'\) are hypotheses of the
clause; the argument uses only the row of \(K\) and its law certificate.
\end{proof}

\begin{proposition}[Closure of residual causality]\label{prop:causality-closure}
Every kind residual (Definition~\ref{def:residual-scheme}), in particular every residual-causality feature, of a description
\(D \in \FATCD{}\) whose content takes one of the recognized forms below,
where a form naming a \(P_i\) projection or unabsorbed \(P_i\)-aligned content
requires the feature to have \(P_i\)-form (paragraph \emph{The \(P_i\)-form}
before Proposition~\ref{prop:probability-closure}),
is licensed under the residual scheme, relative to every well-formed
decomposition record \(\Dec{D}\), at least one of the labels of
\begin{itemize}[topsep=2pt,itemsep=1pt]
  \item \(P_4\)-aligned content, label (i) or (v), as the primary classification when the
    feature has \(P_4\)-form (for instance a directed transition or
    refinement record with its law);
  \item \(P_3\)-aligned content, label (i) or (v), only when the feature has \(P_3\)-form and is specifically a
    route/coherence comparison between generated same-source,
    same-target routes;
  \item \(P_6\)-aligned content, label (i) or (v), when the feature has \(P_6\)-form;
  \item label (i), (ii) or (v), as a member of a composite of \Pfour{} and
    \Psix{}, when directed dependence and provenance audit are carried by a
    finite set of records containing the feature, each of which has
    \(P_4\)-form or \(P_6\)-form (so it satisfies criterion (i) or (v), and also (ii)
    when it contributes required fields to two named active clusters and every
    predicate, comparison, invariant, law or claim field is required for, or an
    input to the check of, either cluster's witness predicate);
  \item an outside-domain structure, when an assertion field of the feature (for instance a
    metaphysical-causation assertion) uses a declared predicate whose extension \(D\) does not record, the asserted property is true in some expansion of \(D\) and false in another, and every other assertion field of that record either also meets the outside-scope condition or is required for \(W_i\) at a cluster aligned with \(P_i\), an input to its check at such a cluster, or named in a seed clause that the record carries;
  \item unabsorbed role-aligned content, label~(v) of
    Definition~\ref{def:residual-scheme}, when it has \(P_i\)-form but
    criterion (i) does not hold for it, for instance because the cluster lacks
    the attachment data of Definition~\ref{def:threshold-attachment}.
\end{itemize}
Classification as \Pfive{} is inadmissible in the absence of quotient
packaging. Each case licenses a label among (i)--(ix) of
Definition~\ref{def:residual-scheme}, so neither (x) nor (xi) is licensed for
the feature. For these recognized forms, the covered content requires no
seventh role; the outside-domain case receives label (viii) without a claim
about roles beyond the covered scope. None of these classifications uses the
recognition hypothesis (Definition~\ref{def:recognition-hypothesis}); content
of none of these forms is governed by it. An interventional mechanism is covered
when recorded as a law-bearing transition or refinement record of \(P_4\)-form, for instance one whose only assertion field is its law certificate (first bullet); a further assertion field serving no witness can leave it licensed only (x); a
counterfactual assertion is covered only when it takes one of the listed forms,
and otherwise falls under the recognition hypothesis.
\end{proposition}

\begin{proof}
Let \(r\) be a residual-causality feature (a kind residual) of
\(D \in \FATCD{}\), causal content being taken from the hierarchy of~\citet{Pearl2009}, and pass
\(r\) through the residual scheme of Definition~\ref{def:residual-scheme} under
Proposition~\ref{prop:alignment-rules}. The forms of the proposition cover
interventional mechanisms recorded as transitions of \(P_4\)-form; counterfactual assertions
not of these forms fall under the recognition hypothesis. A causal feature presents a
directed dependence, so we case on how that directedness is carried.
Role-aligned and composite forms are labelled as in the paragraph
\emph{The \(P_i\)-form} before Proposition~\ref{prop:probability-closure}: each
record receives (i) or (v), and (ii) under the two-cluster condition stated
there.

If \(r\) has \(P_4\)-form as a directed transition, refinement, or
staged-dependence law, it receives label (i) or (v) by the paragraph
\emph{The \(P_i\)-form}; this is the primary classification. The record
kind and its law alone do not establish the assertion-field guard
defining \(P_4\)-form. One case in which directedness routes elsewhere is when \(r\) is
specifically a route/coherence comparison between generated
same-source, same-target routes carrying a non-reducibility certificate of
either alternative in the \(W_3\) row; only then does it supply the \Pthree{}
witness, have \(P_3\)-form and receive label (i) or (v). We must not assign a
generic causal arrow to \Pthree{}, since by the alignment rules a
generic transition is not a route comparison of the \Pthree{} type. If
\(r\) is instead a provenance or dependency audit --- event reference
with a provenance or trace record --- it has \(P_6\)-form. A \(P_6\)-form
record receives (i) when its named active cluster carries the certified trace
check and honest audit required by \(W_6\); a covered seed without that full
witness receives (v). When directed dependence and provenance audit occur together, take
their finite record set \(R_f\). Its \(P_4\)-form and \(P_6\)-form
members are licensed (i) or (v) individually by their \(P_i\)-form; together they have
the stated composite reading.
Finally, if \(r\) records a metaphysical-causation assertion whose declared
predicate has no recorded extension in \(D\), its truth is not certified by
the finite checks of \(D\). In the outside-domain case, the asserted property is true in some
expansion of \(D\) and false in another. Together with the stated
condition on the remaining assertion fields, this satisfies
criterion (viii).

Classification as a \Pfive{} projection is not available here: by
Proposition~\ref{prop:alignment-rules} a \Pfive{} projection requires a
quotient relation and packaging map, and a causal feature carrying no
quotient packaging supplies none; the directedness of a causal chain is
not packaging data. Hence the \Pfive{}-inadmissibility caveat holds.
The cases above are the recognized forms of residual-causality content; a
causality feature whose content is of none of these forms is not classified here
and falls under the recognition hypothesis
(Definition~\ref{def:recognition-hypothesis}). The closure is scoped and makes no
scope-free claim.
\end{proof}

\begin{proposition}[Closure of residual agency]\label{prop:agency-closure}
Every kind residual (Definition~\ref{def:residual-scheme}), in particular every residual-agency feature, of a description
\(D \in \FATCD{}\) whose content takes one of the recognized forms below,
where a form naming a \(P_i\) projection or unabsorbed \(P_i\)-aligned content
requires the feature to have \(P_i\)-form (paragraph \emph{The \(P_i\)-form}
before Proposition~\ref{prop:probability-closure}),
is licensed under the residual scheme, relative to every well-formed
decomposition record \(\Dec{D}\), at least one of the labels of
\begin{itemize}[topsep=2pt,itemsep=1pt]
  \item label (i), as a constituent of an active \Ptwo{}, \Pfour{} or \Psix{}
    projection, when the feature has the
    corresponding \(P_i\)-form and \(\Dec{D}\) names, for an active \(P_i\) projection, the cluster \(C\) of the paragraph \emph{The \(P_i\)-form} (for the seed branch of that paragraph, any named cluster containing the feature);
  \item label (i), (ii) or (v), as a member of a composite of \Ptwo{},
    \Pfour{}, and \Psix{}, when the feature is one of an objective or
    specification record of \(P_2\)-form, a
    transition, regime, or staged selection record of \(P_4\)-form, and an
    audited decision-provenance record of \(P_6\)-form, each licensed (i) or (v) by its
    \(P_i\)-form;
  \item an empirical-bridge annotation, when the feature is itself an empirical-bridge annotation identifying an admissible bridge claim (Definitions~\ref{prop:honest-bookkeeping} and~\ref{prop:empirical-bridge}) whose \(\Audit{D}\) entry carries the five items of Definition~\ref{prop:empirical-bridge}, with no fields beyond its typing, targets and that identification;
  \item an outside-domain structure, when a proposed irreducible intentional
    or teleological predicate has no recorded extension in \(D\), the asserted property is true in some expansion of \(D\) and false in another, and every other assertion field of that record either also meets the outside-scope condition or is required for \(W_i\) at a cluster aligned with \(P_i\), an input to its check at such a cluster, or named in a seed clause that the record carries;
  \item unabsorbed role-aligned content, label~(v) of
    Definition~\ref{def:residual-scheme}, when it has \(P_i\)-form but
    criterion (i) does not hold for it, for instance because the cluster lacks
    the attachment data of Definition~\ref{def:threshold-attachment}.
\end{itemize}
A content record targeted by such an empirical-bridge annotation is not thereby classified; it is classified by its own fields and is licensed label (ix-c) exactly when either branch of that criterion holds.
An update map receives (i) when it lies in a named active \(W_1\) cluster and
its assertion fields meet criterion (i); otherwise it receives (v) when it
carries a covered \Pone{} seed whose clause names each of its predicate,
comparison, invariant, law or claim fields, or belongs to a failed \(W_1\)
offer each of whose such fields is an input to the recorded check; and (ix-c)
when it supplies no witness field and has no such field; an update map of none of these three forms must be tested against the other recognized forms of this proposition; only when none applies is it not classified here and left to the recognition hypothesis. Each case licenses a label among (i)--(ix) of
Definition~\ref{def:residual-scheme}, so neither (x) nor (xi) is licensed for
the feature. For these recognized forms, the
covered content requires no seventh role; the outside-domain case receives
label (viii) without a claim about roles beyond the covered scope. None of
these classifications uses the recognition hypothesis
(Definition~\ref{def:recognition-hypothesis}); content of none of these forms
is governed by it.
\end{proposition}

\begin{proof}
Let \(r\) be a residual-agency feature (a kind residual) of
\(D \in \FATCD{}\), agency being taken in the planning sense
of~\citet{Bratman1987}, and pass it through the residual scheme of
Definition~\ref{def:residual-scheme} under
Proposition~\ref{prop:alignment-rules}. An agency feature in admissible
form presents an objective pursued through selection and recorded for
later inspection; we case on whether that content reduces to role
witnesses or asserts more. Role-aligned and composite forms are labelled as in the paragraph
\emph{The \(P_i\)-form} before Proposition~\ref{prop:probability-closure}: each
record receives (i) or (v), and (ii) under the two-cluster condition stated
there.

If \(r\) has just one of the stated \(P_2\)-, \(P_4\)- or \(P_6\)-forms and
receives (i), criterion (i) supplies its named active cluster, \(W_i\)
witness and attachments; hence it is the corresponding single-role
projection. In the composite case, take the finite record set \(R_f\) carrying
the objective or specification, staged selection, and decision
provenance. By the stated \(P_2\)-, \(P_4\)- and \(P_6\)-form
hypotheses, each member receives (i), (ii) or (v) individually.
Together those records have the asserted composite reading. If instead \(r\) is an empirical-bridge annotation whose bridge claim is admissible and has an
audit entry with the five items of Definition~\ref{prop:empirical-bridge}, and with
no fields beyond its typing, targets and that identification, it
receives (vii). A content record targeted by such an annotation is not thereby classified; it is classified by its own fields and is licensed label (ix-c) exactly when either branch of that criterion holds. If \(r\) records a
proposed irreducible intentional or teleological predicate whose extension
\(D\) does not record, its truth cannot be checked over \(D\). In the outside-domain case, the asserted property is true in some
expansion of \(D\) and false in another. Together with the stated
condition on the remaining assertion fields, this satisfies
criterion (viii). A finite audited intentional-content claim with a recorded
predicate extension and no recognized six-role disposition remains subject to
the recognition hypothesis.

Classification as a \Pone{} projection is not available here merely
because agency selects an update: by
Proposition~\ref{prop:alignment-rules} a \Pone{} projection requires a
selected update together with source and target relation records of declared use quotient and an
explicit descent-compatibility, descent-obstruction, or
update-vs-quotient mismatch witness. A selection event carrying no such
witness supplies no \Pone{} witness. Its update map receives (i) when it lies
in a named active \(W_1\) cluster and its assertion fields meet criterion (i);
otherwise it receives (v) when it carries a covered \Pone{} seed whose clause
names each of its predicate, comparison, invariant, law or claim fields, or
belongs to a failed \(W_1\) offer each of whose such fields is an input to the
recorded check; and (ix-c) when it supplies no witness field and has no such
field. It has \(P_4\)-form only when recorded as a
law-bearing transition record. This is the stated \Pone{} caveat. The cases
above are the recognized forms of residual-agency content; content of none of
these forms falls under the recognition hypothesis
(Definition~\ref{def:recognition-hypothesis}). The closure is scoped and asserts
nothing beyond it.
\end{proof}

\begin{remark}[Scope of the closures]\label{rem:threat-closures-scope}
Propositions~\ref{prop:probability-closure}--\ref{prop:agency-closure}
close the three threats within the covered scope; they do not make
scope-free claims. A feature that falls outside the covered scope is
classified as an outside-domain structure via
Definition~\ref{def:residual-scheme}, not absorbed into a primitive role
by default. Each named form is checked against its label criterion directly, without
requiring recognition of every feature of \(D\); the propositions give the
applicable label criterion for each named form and state when \Pfive{} and
\Pone{} projections are unavailable. Each recognized form in
Propositions~\ref{prop:probability-closure}--\ref{prop:agency-closure} is
alternative (R1), (R2), (R3) or (R5) of Definition~\ref{def:recognition-hypothesis}:
each individual record in a recognized form satisfies (R1), (R2), (R3) or (R5), and
composite forms are assessed member by member; the propositions verify those alternatives for
the listed forms, and add that a \Pfive{} reading (causality) or a \Pone{}
reading (agency) is unavailable without quotient-packaging or descent data.
\end{remark}

\subsection{The recognition hypothesis and the classification theorem}\label{subsec:classification-theorem}

Each alternative of Definition~\ref{def:recognition-hypothesis} is the
decomposition-independent form of label criteria of
Definition~\ref{def:residual-scheme}: (R1) of (i) or (v); (R2) of (v), (vi),
(vii), (ix-a) and (ix-b); (R3) of (ix-c); (R4) of (iii) and (iv); and (R5) of
(viii). Theorem~\ref{thm:upper-bound-classification} verifies that this
matching does not depend on the chosen decomposition record.

The hypothesis below is imposed on three sets: the kind residuals of Definition~\ref{def:residual-scheme} (content records and structural annotations, whose role reading is not fixed by their kind), all bookkeeping annotations, and the set \(\Residual{D}\) of each well-formed decomposition record (Definition~\ref{def:decomposition-record}). Stated for every record, it
does not depend on a choice among them. Since \(\Residual{D}\) consists of kind
residuals and bookkeeping annotations, the third clause below follows from the first two: \(D\) satisfies the recognition
hypothesis if and only if every kind residual and every bookkeeping annotation of
\(D\) is recognized, a condition that refers to no decomposition
record.

\begin{definition}[Recognition hypothesis]\label{def:recognition-hypothesis}
Terms used below: \emph{aligned}, \emph{contributes a required field},
\emph{required for}, \emph{input to the check}, \emph{declared operation},
\emph{independent assertion} and \emph{outside the covered scope} are defined in
the paragraph \emph{Terms used in the label criteria} after
Definition~\ref{def:residual-scheme}; \emph{covered seed}, \emph{named seed
field} and \emph{failed offer} in Definition~\ref{def:channel-status};
\emph{honestly recorded} (of an offer: its audit is honest,
Definition~\ref{def:audited}(b),(d)).
The \emph{closure content} of a feature \(r\) is the finite induced
subdescription consisting of \(r\), its fields, and all declared records
reached by following its references in \(D\) or read by its formula-valued fields; it is used in
Remark~\ref{rem:recognition-open} (paragraph \emph{Verification}). The alternatives below are evaluated in the
whole of \(D\), and a reference makes data available for checking without by
itself certifying a role witness. A description \(D \in \FATCD{}\)
satisfies the \emph{recognition hypothesis} when every kind residual \(r\)
of \(D\) (Definition~\ref{def:residual-scheme}), every bookkeeping annotation \(r\) of \(D\) (whether or not a decomposition record assigns it to a channel), and every
element \(r\) of \(\Residual{D}\) for every well-formed decomposition record
of \(D\) (Definition~\ref{def:decomposition-record}), is \emph{recognized}, that is,
satisfies at least one of the following:
\begin{enumerate}[label=(R\arabic*),topsep=2pt,itemsep=1pt]
  \item (first branch) \(r\) itself contributes a required field to a finite cluster aligned
    with one of the six witness types of the \emph{Witness types} table of Section~\ref{sec:primitive-role-architecture},
    or itself carries a covered seed for one of them
    (Definition~\ref{def:channel-status}), and every predicate, comparison, invariant, law or claim field \(\varphi\) of \(r\) is, for some index \(i_\varphi\) and some finite cluster \(C_\varphi\ni r\) aligned with \(P_{i_\varphi}\), required for \(W_{i_\varphi}\) at \(C_\varphi\) or an input to its check, or is named in a seed clause that \(r\) carries (so \(r\) carries no assertion beyond these fields; content beyond them must be recorded in a separate declared content record, which is classified on its own); or (failed-offer branch)
    \(r\) is a member of the cluster of an honestly recorded failed offer for
    some \(P_i\), and each of its predicate, comparison, invariant, law or claim
    fields is an input to the recorded check of \(W_i\) on that cluster; merely
    referencing a witness cluster does not suffice;
  \item \emph{Structural annotations.} \(r\) is one of the six structural annotations of
    Definition~\ref{def:annotations}; its fields contain only its declared
    operation, targets, and the finite typing, checking, preserved or
    suppressed data, losses and nonclaims required for that operation, and it
    makes no independent assertion (both terms as defined after
    Definition~\ref{def:residual-scheme}); in particular an instrument-visibility
    annotation whose fields other than its typing and target fields are those
    required by Definition~\ref{prop:visibility}, a lift annotation whose fields are those
    required by Definition~\ref{prop:forgetting-lifts}, and an
    empirical-bridge annotation meeting criterion (vii), whose fields are only
    its typing, targets and the identification of its bridge claim and five
    items, satisfy this
    alternative, provided each of their formula-valued fields meets the
    condition of the next sentence. A formula-valued field of a structural annotation (a threshold predicate, a preserved or lost invariant) is part of its declared operation only when it reads at least one record, every record whose values it reads is listed in the annotation's target field, and each such record satisfies (R1); otherwise it is an independent assertion, and this alternative fails.
    \emph{Bookkeeping annotations.} The alternative is also met when \(r\) is one of the seven bookkeeping annotations (for \(D\in\FATCD{}\) every condition below except the last, on field roles, holds automatically by Definitions~\ref{def:finite}--\ref{def:fatcd}) of
    Definition~\ref{def:annotations}, it is finitely represented, has finite typed fields, and every reference it carries resolves to a declared finite record of \(D\) (Definitions~\ref{def:finite} and~\ref{def:closed}), it is honest if it is an audit annotation, it carries the fields Definition~\ref{def:audited}(g) requires of its kind, and, if it is a claim annotation, it is referred to by an honest audit annotation and satisfies the representation condition (Definition~\ref{def:audited}(e)), and every field of
    \(r\) other than the statement and hypothesis fields of a claim annotation,
    the certificate field of an audit annotation and its conclusion field when its conclusion is true in every well-typed interpretation or is represented, in the sense of Definition~\ref{def:annotations}, by fields of the content records named by the claim it audits, and the
    formula schema of a witness-schema annotation has field role datum or
    reference; a substantive assertion in a claim statement is represented by
    content records (Definition~\ref{def:annotations}), which are assessed
    under (R1)--(R5);
  \item \(r\) is role-neutral data without an assertion of its own: \(r\)
    supplies no witness field and has no predicate, comparison, invariant,
    law or claim field. Typical instances are the inputs that a specified field
    of a feature satisfying (R1) uses in checking a required or seed field, and
    declared objects or states that no role datum uses; also a predicate record supplying no witness field whose only assertion field is a formula with at least one free variable whose truth value is not the same under all assignments respecting the declared sorts of its free variables, using only predicates whose extensions \(D\) records, and quantifying only over sorts whose carriers \(D\) records, each free variable ranging over a record kind or the rational sort, referenced by no claim annotation, and referenced by certificates only through the abbreviation \(s(t)\) of the paragraph \emph{Check language} (a definition, which has no truth value over \(D\));
  \item \(r\) begins a finite chain \(r=r_0,\ldots,r_n\) of distinct
    declared records, with \(n\geq1\), such that, for each \(j<n\), a field of \(r_j\) records a checked
    equality or isomorphism of presentation data with specified fields of
    \(r_{j+1}\), or a checked level translation to \(r_{j+1}\); and, for each \(j<n\), apart from one comparison field recording the equality, isomorphism or level translation with \(r_{j+1}\), and the datum and reference fields it compares, \(r_j\) has no predicate, comparison, invariant, law or claim field. The terminal \(r_n\) satisfies (R1), (R2), (R3) or
    (R5). An equality or isomorphism of presentation data (a value-preserving
    relabelling of datum fields; no structural isomorphism is meant) is a finite
    bijection \(\beta\), recorded in the comparison field, from specified datum fields that are not term-valued
    (or the member list of an object record, a reference field) of \(r_j\) to datum fields that are not term-valued (or the member list of an object record, a reference field) of \(r_{j+1}\), with a
    correct certificate, recorded true, that each recorded value equals that of its \(\beta\)-image;
    its formula is a conjunction of equalities between record identifiers,
    recorded field values and entries of recorded lists (written with indexing,
    as in the running example of Section~\ref{sec:fatcd-domain}). A checked level translation from \(r_j\) to \(r_{j+1}\) is a
    comparison field of \(r_j\) referencing a translation annotation carried
    in \(\Audit{D}\) whose source and target contexts are levels recorded by level annotations whose target fields list \(r_j\) and \(r_{j+1}\), respectively, together with a finite
    bijection \(\beta\) of specified datum fields as above and a correct certificate, recorded true,
    that is such a conjunction of equalities,
    that each value equals its \(\beta\)-image;
  \item \(r\) lies outside the covered scope
    (Definition~\ref{def:residual-scheme}): a field of \(r\) asserts a property whose check is not finite over the records of \(D\): it quantifies over a declared domain for which \(D\) supplies no finite checking range, or uses a declared predicate whose extension \(D\) does not record, and its truth value is not fixed by the recorded values of \(D\) (it is true in some expansion of \(D\) and false in another, in the sense of the paragraph \emph{Certificates, assertion language and laws} of Section~\ref{sec:primitive-role-architecture}; for a formula with free variables: for some assignment \(\alpha\) of its free variables and two expansions of \(D\) whose carriers contain the values \(\alpha\) gives to variables of unlisted sorts, it is true under \(\alpha\) in one and false under \(\alpha\) in the other). In addition, every other predicate, comparison, invariant, law or claim field of \(r\) either has this property or is required for \(W_i\) at a cluster \(C\) aligned with \(P_i\), an input to the check of \(W_i\) at such a cluster, or named in a seed clause that \(r\) carries; \(r\) is then
    recognized as outside-domain structure.
\end{enumerate}
For \(D\in\FATCD{}\), every bookkeeping annotation whose fields other than the
exempt fields listed in the second clause of (R2) are datum or reference fields
satisfies the
second clause of (R2); hence, when all bookkeeping annotations of \(D\) have this
form, \(D\) satisfies the recognition hypothesis if and only if every
kind residual of \(D\) is recognized.
A feature satisfying (R1) is absorbed under label (i) when a decomposition record
places it in an active projection under
Proposition~\ref{prop:alignment-rules}, and is covered under (v) when it
remains unabsorbed. Remark~\ref{rem:recognition-open} shows that some
descriptions in \FATCD{} do not satisfy this hypothesis.
\end{definition}

Alternatives (R1)--(R5) are record-independent forms of label criteria: (R1)
of (i) or (v); (R2) of (v), (vi), (vii), (ix-a) or (ix-b); (R3) of (ix-c);
(R4) of (iii) or (iv); and (R5) of (viii) (the table in the proof of
Theorem~\ref{thm:upper-bound-classification}).

\begin{remark}[The recognition hypothesis is a genuine restriction]\label{rem:recognition-open}
The recognition hypothesis is an explicit assumption on \(D\), not a consequence
of the \FATCD{} definitions, and some descriptions in \FATCD{} do not satisfy
it. The remark establishes five facts, each under the paragraph named:
\(D_{\mathrm{norm}}\in\FATCD{}\setminus\FATCD_{\mathrm{RH}}\), so
\(\ScopedExactSix{\FATCD{}}\) fails (\emph{The counterexample},
\emph{Verification}); a transition-row presentation of the same weights lies
in \(\FATCD_{\mathrm{RH}}\) (\emph{Repair by a transition row}); every closed
check-language formula without atomic \(W_j\) formulas has a covered
\Ptwo{}-seed presentation, given two distinct declared records of one kind (or after adjoining them), whenever the completed description is again in \FATCD{} (\emph{Wrapping a closed formula as a \Ptwo{} seed}; membership can fail when a correct certificate, or a realization relation's target formula, quantifies over all declared records of an adjoined kind);
membership in \(\FATCD_{\mathrm{RH}}\) depends on presentation (\emph{Scope
depends on presentation}); and label (xi) is realized by \(D_{\mathrm{scr}}\)
(\emph{A description passing the screen}). The paragraphs \emph{One record suffices for recognition}, \emph{Other unrecognized assertions} and \emph{Open question} add a one-record failed-offer presentation, further unrecognized forms with the (R2) and (R3) presentation devices, and the open question. A record without an assertion of its own --- one with no witness field and
no predicate, comparison, invariant, law or claim field, such as a declared
object record --- is recognized by (R3). The normalization predicate in the
following description is not recognized.
\par\emph{The counterexample \(D_{\mathrm{norm}}\).} Let \(D_{\mathrm{norm}}\), the description of Remark~\ref{rem:fatcd-minimality}, consist of state records \(w_1,\dots,w_n\)
(\(n\ge1\)), each with a nonnegative rational value field \(v(w_k)\), the values
summing to one (for instance \(n=2\) and \(v(w_1)=v(w_2)=1/2\)); a predicate
record \(N\), declared as a macro-specification, whose formula asserts
\(\bigwedge_k v(w_k)\ge0\wedge\sum_k v(w_k)=1\); a claim annotation naming \(N\) whose statement is the
formula of \(N\); and an honest audit annotation listing the certificate
\((\bigwedge_k v(w_k)\ge0\wedge\sum_k v(w_k)=1,\mathrm{true})\) and certifying that formula.

\emph{Verification.} \(D_{\mathrm{norm}}\) is finite, closed and
audited, so \(D_{\mathrm{norm}}\in\FATCD{}\). It contains no map, relation,
comparison, path, derivation, order, transition, refinement, event,
provenance, trace, replay or check record; in particular the closure content
of \(N\), namely \(N\) with \(w_1,\dots,w_n\), contains none. So no cluster
satisfies any \(W_i\)
and no record is a covered seed (Definition~\ref{def:channel-status}). Hence
\(N\) does not satisfy (R1); being a content record, not an annotation, it does not satisfy (R2); its predicate field contains a closed formula and it is referenced by a claim annotation, so neither branch of (R3) applies;
it records no presentation or level translation, so not (R4); and \(N\) gives
its predicate by the explicit formula \(\bigwedge_k v(w_k)\ge0\wedge\sum_k v(w_k)=1\) over recorded field
values; this formula uses no declared predicate symbol and quantifies over no
sort, so its check is finite over \(D_{\mathrm{norm}}\) and (R5) fails.
Although \(N\) is declared a macro-specification, it is not a covered seed:
\(D_{\mathrm{norm}}\) has no realization relation to a declared nonempty
candidate class, and no realization relation can target \(N\), since
Definition~\ref{def:raw-record} requires its target formula to have exactly one
free variable. For the same reasons no
criterion (i)--(ix) of Definition~\ref{def:residual-scheme} holds for \(N\):
it is not in any cluster or composite, re-presents or restates nothing, is not
an annotation, relates no instrument or substrate observable, is inside the
covered scope, and carries an assertion. With no failed offer and no outside-scope declaration,
every channel of \(D_{\mathrm{norm}}\) is \(\absentrec\), so no active projection absorbs \(N\), and \(N\) receives
label (x) in every well-formed decomposition record. So \(\ScopedExactSix{\FATCD{}}\)
fails, and the six roles are not \emph{recognition-complete} for \FATCD{}: some in-scope
feature of some description in \FATCD{} is placed by none of (i)--(ix).
Label (x) records an unresolved candidate, not a seventh role: (xi) would
further require the recorded witness of category (xi), itself a recorded
screen rather than the identification of a new closure role.

\emph{Repair by a transition row.} For distribution-valued weights the hypothesis constrains presentation, not content. For nonnegative
rational weights summing to one, replace \(N\) by a transition record
\(K:\{*\}\to\mathrm{Dist}(\{w_1,\ldots,w_n\})\) whose unique row is the list
of terms \(v(w_1),\dots,v(w_n)\), so that \(K\) reads the value fields of the
\(w_k\) and records no copies. Declare a state record \(*\), a one-point object record \(\{*\}\) listing it, and an
object record \(X\) listing \(w_1,\dots,w_n\), give \(K\) a codomain reference to
\(X\) and a datum field marking it distribution-valued, and take as the law certificate
of \(K\) the pair \((\bigwedge_k v(w_k)\ge0\wedge\sum_k v(w_k)=1,\mathrm{true})\),
which expresses that the row lies in \(\mathrm{Dist}(X)\). Retarget the claim and its audit to \(K\), replacing every reference
to \(N\) by a reference to \(K\). Set the claim's statement and the
audit's conclusion to the law-certificate formula, and replace the
audit's certificate list by that certificate, recomputed true,
naming \(K\) and its declared inputs. Thus \(K\) represents the
statement (Definition~\ref{def:annotations}), and the audit is honest. Then \(K\) is a covered \Pfour{} seed
satisfying (R1); \(*\), \(\{*\}\), \(X\) and the \(w_k\) carry no assertion field and
satisfy (R1) or (R3); the claim and audit annotations satisfy (R2); and the
resulting description lies in \(\FATCD_{\mathrm{RH}}\).

\emph{Wrapping a closed formula as a \Ptwo{} seed.} In outline: the construction adjoins a predicate record \(s\) with one free variable over two records \(c\ne c'\), a realization relation recording \(s(c)=\mathrm{true}\) and \(s(c')=\mathrm{false}\), and a one-class packaging map \(q\) with a comparison record; it chooses between two formulas for \(s\) according to the value of \(\varphi\) in the completed description \(D'\). More generally, let \(\varphi\) be any closed check-language formula containing
no atomic \(W_j\) formula and
\(c\ne c'\) declared records of one kind. Let \(D'\) be the final description, obtained once every record of this construction (including the records \(X_c,*,Q,q\) and the comparison record introduced under \emph{Activation} below, the \Ptwo{} claim and its audit, all attachment annotations and any fresh candidate records, when used) is adjoined and any removal below is made. The predicate record is \(s(r):\equiv\varphi\wedge r=c\) if \(\varphi\) holds in \(D'\), and \(s(r):\equiv\varphi\vee r=c\) otherwise. \par\emph{Choice of branch.} Fix all additions, removals and retargetings, and all non-formula field values, before choosing the branch. The two branches differ only in the formula field of \(s\), the retargeted claim's statement and its audit's conclusion. A check-language formula never reads a statement or conclusion (formula-valued fields are not terms) and reads a predicate record's formula only through an abbreviation \(s(t)\) naming it; \(s\) is adjoined, so \(\varphi\) contains none. Hence \(\varphi\) has the same value in both candidates for \(D'\). Its value agrees with its value in \(D\) provided every field it reads is unchanged and every quantifier has the same checking range. The construction applies when \(D'\) is again in \FATCD{}. Adjoining records can falsify a correct certificate of \(D\) whose quantifier ranges over all declared records or over the declared records of an adjoined kind (for instance \((\forall x\in\mathrm{State}\,(x=w_1\vee x=w_2),\mathrm{true})\), which the adjoined state record \(*\) falsifies). A sufficient condition is that every retained realization relation still lists exactly its satisfying members and records correct evaluations, every retained certificate and law remains correct, every retained audit remains honest, and all additions, removals and retargetings preserve typing, reference closure, representation and the required annotation fields. \par\emph{Seed and claim.} This predicate record, declared a macro-specification whose variable \(r\) is declared to range over the common kind of \(c\) and \(c'\), with
a realization relation declaring \(K=\{c,c'\}\), listing exactly \(c\) as
satisfying, and recording both evaluations, is a covered \Ptwo{} seed. Let the
realization relation carry the evaluations as certificates
\((s(c),\mathrm{true})\) and \((s(c'),\mathrm{false})\). The record that carried \(\varphi\) as an assertion field, if any (such as
\(N\)), may be removed only if neither \(\varphi\) nor its named inputs
reference it, it is neither \(c\) nor \(c'\), and every incoming reference is a
bookkeeping reference that can be retargeted to \(s\) and the realization
relation while preserving typing, representation and audit honesty; otherwise
that record is retained and classified separately; a claim whose statement
is \(\varphi\) is retargeted to \(s\) and the realization relation, with statement
\(\varphi\) in the first branch and \(\neg\varphi\) in the second; its audit lists
both certificates and concludes that statement. The statement is entailed by
the signed certificates, so the claim satisfies the representation condition
of Definition~\ref{def:annotations} and the audit is honest. Its
predicate field is a named seed field, so it satisfies (R1). It receives (i)
when criterion (i) holds and (v) otherwise. \par\emph{Activation.} Adjoining an object record \(X_c\)
listing \(c,c'\), a state record \(*\) with a quotient object record \(Q\)
listing it, a map record \(q:X_c\to Q\) of declared use packaging with
\(q(c)=q(c')=*\), and a comparison record carrying
\((q(c)=q(c')\wedge s(c)\wedge\neg s(c'),\mathrm{true})\) gives the following. Let \(C\) contain \(c,c',s\), the realization
relation, \(X_c\), \(*\), \(Q\), \(q\) and the comparison record. Then
\(W_2(C)\) holds: the realization relation records both evaluations on
\(K=\{c,c'\}\), and the single kernel class of \(q\) meets the realizer set
\(\{c\}\) and its complement \(\{c'\}\). Record a claim naming \Ptwo{} and \(C\)
with statement \(W_2(C)\), and an honest audit listing the row's content
certificates and \((W_2(C),\mathrm{true})\) and concluding \(W_2(C)\), and
supply all attachments and level and lift conditions of
Definition~\ref{def:role-projection}; fresh candidate records may be used to
avoid pre-existing level, lift or forgetting annotations meeting \(C\). The
resulting projection is active. Since \(q\) is a covered \Pfive{} seed, in \(D'\) channel \Ptwo{} is active and channel \Pfive{} is not \(\absentrec\); the transition-row repair likewise makes \Pfour{} at least \(\belowthr\). Removing \(s\) from this cluster falsifies
\(W_2\), and its sole predicate field is required for that witness, so
criterion (i) applies in a decomposition record naming this projection. \par\emph{What is recorded.} This construction records the truth value of \(\varphi\) in \(D'\) in recognized
seed content: in the first branch the recorded evaluation \(s(c)=\mathrm{true}\)
entails \(\varphi\), while in the second branch \(s(c')=\mathrm{false}\) entails
\(\neg\varphi\). Thus an assertion true in \(D'\) is recorded with its checked content; an assertion false in \(D'\) is represented by its checked negation. Agreement with its truth value in \(D\) requires the preservation conditions above. The transition-row
presentation above likewise preserves the checked normalization content for
nonnegative weights of pairwise distinct records.

\emph{One record suffices for recognition.} Let \(c\) be a declared record, \(s(r):\equiv\varphi\wedge r=c\), and let a realization relation with candidate class \(K=\{c\}\) and target \(s\) carry \((s(c),\mathrm{true})\) if \(\varphi\) holds in the completed description and \((\neg s(c),\mathrm{true})\) otherwise. Fix all additions and fields other than that certificate's formula and the relation's list of satisfying members before selecting the branch. Assume \(\varphi\) reads neither changing field in either completion, including through quantified field lookups or predicate abbreviations. Then \(\varphi\) has the same value in both completions, so the branch choice is well defined. The construction applies when the completed description, with \(s\), the relation, the claim and its audit adjoined, is again in \FATCD{}. Record a claim naming \Ptwo{} and \(C=\{c,s,\mathrm{Real}\}\) with an honest audit listing that certificate and \((W_2(C),\mathrm{false})\). Since \(K=\{c\}\), no kernel class can meet both the realizer set and its complement in \(K\), and no evaluation outside \(K\) is recorded in \(C\); hence \(W_2(C)\) is false in every well-typed interpretation; the formula of \(s\) and the certificate are inputs to the recorded check, so \(s\) and the relation satisfy the failed-offer branch of (R1) and receive (v-c), and the certificate entails \(\varphi\) or \(\neg\varphi\). Channel \Ptwo{} is then \(\collapsedbc\) unless active. The two-record proviso is needed only for the covered-seed and active presentations.

\emph{Scope depends on presentation.} The covered scope also depends on presentation. Replace the formula of \(N\)
by \(\mathrm{IsDist}(v(w_1),\ldots,v(w_n))\), declare \(\mathrm{IsDist}\) by a
predicate record of arity \(n\), each argument place declared rational-sorted, with no formula field and give no relation
record for it, so that \(D\) does not record its extension (the declaring
record has no assertion field and satisfies (R3); had the declaring record
given an explicit formula in its \(n\) rational argument variables, such as \(\bigwedge_k x_k\ge0\wedge x_1+\dots+x_n=1\), \(D\) would record the extension and \(N\) would
again receive (x)), set the claim's statement
to \(\mathrm{IsDist}(v(w_1),\ldots,v(w_n))\), and let the audit list no certificate, conclude the valid formula
\(v(w_1)=v(w_1)\), and list the \(\mathrm{IsDist}\) field of \(N\) in its
non-certification field. This description remains honestly
audited. Its records \(w_k\) and the declaring record of \(\mathrm{IsDist}\)
satisfy (R3), its claim and audit annotations satisfy the second clause of
(R2), and \(N\) satisfies (R5) and receives (viii) rather than (x); so the
description lies in \(\FATCD_{\mathrm{RH}}\) while \(D_{\mathrm{norm}}\) does
not, and no role is certified for \(N\). Hence membership in \(\FATCD_{\mathrm{RH}}\) is not
invariant under replacing an explicit formula by a declared predicate name whose extension \(D\) does not record,
and Theorem~\ref{thm:scoped-exact-six} asserts nothing about content presented
that way. More generally, replace at least one assertion field of an unrecognized
kind residual by a declared predicate name whose extension the resulting
description does not record. The changed record satisfies (R5) provided every
other assertion field, evaluated in the resulting description, also lies
outside the covered scope, is required for or an input to \(W_i\) at an
aligned cluster containing the record, or is named in a seed clause the record
carries. Unchanged records must be recognized separately. The result lies in
\(\FATCD_{\mathrm{RH}}\) if it is again in \FATCD{} (in particular, no replaced field is the formula of a realization relation's target),
every remaining kind residual is recognized, and every bookkeeping annotation
has only datum and reference fields besides the exempt fields of (R2).
Theorem~\ref{thm:scoped-exact-six} certifies no role for the replaced
fields.

\emph{Other unrecognized assertions.} A standalone predicate record stating a conservation or balance law, or an
assembled obstruction or holonomy, by a closed check-language formula over
recorded inputs is not recognized as presented if it carries no covered seed,
contributes no required witness field, is not an input to a prescribed witness
check, and records no checked presentation or level translation. Under these
conditions (R1)--(R5) fail as for \(N\).
An in-scope content record that contributes a required field to an aligned
cluster or carries a covered seed, and whose assertion fields are each required
at, or inputs to the check at, some aligned cluster containing it, possibly
different clusters for different fields, satisfies (R1); one with no assertion
field and no witness field satisfies (R3).
Label (x) is also licensed for an unclaimed predicate record whose closed
formula is true or false in \(D\), such as \(0+0=0\) or \(0+0=1\) between rational constants: label (x) marks an assertion placed
by none of (i)--(ix), which need be neither true nor a candidate role. By contrast, the same record
with formula \(\theta\wedge r=r\) for a fresh variable \(r\), still unclaimed and
referenced by no certificate, does not satisfy the second branch of (R3), since its value does not depend on
\(r\), and still receives (x). Conversely, if \(c\) is a declared record of the kind over which \(r\) ranges and that kind has a second declared record, the unclaimed record with formula \(\theta\wedge r=c\) (for \(\theta\) true in \(D\)) or \(\theta\vee r=c\) (for \(\theta\) false in \(D\)), referenced by no certificate, satisfies the second branch of (R3), since its value at \(c\) differs from its value at the other records of that kind; for unclaimed closed assertions this branch constrains presentation, not content. The same holds for the structural-annotation branch of (R2): for a closed check-language formula \(\theta\) that reads no record values outside a binary relation record \(E\) satisfying (R1), and a fixed declared record identifier \(u\), a threshold annotation targeting \(\{E\}\) with predicate \(\theta\wedge(E(u,u)\vee\neg E(u,u))\) satisfies (R2) and receives (ix-b) (no cluster \(\{E\}\) satisfies any \(W_i\), so (v-b) cannot apply), whether or not \(\Audit{D}\) records the predicate checked true, provided the description with this annotation adjoined is again in \FATCD{}. Since the
\(W_6\) certificate states a replay or provenance property of its trace, an
arbitrary assertion about recorded values does not become a required
\(W_6\) field by adjoining a trace.

\emph{A description passing the screen.} Let \(D_{\mathrm{scr}}\) consist of
index records \(0,1\), an object record \(X\) listing them, an order record \(o\) for
\(0<1\) with its law certificate, a relation record \(r\) (the feature to be screened; not the free variable \(r\) of a macro-specification) of declared use
other carrying a finite table field \(c\) with \(c(0)=c(1)=1\) and a datum field
holding the list \((\mathrm{id}_X)\), with
certificates of extensivity, monotonicity, \(c\circ c=c\), closure of the list
under composition and \(c\ne\mathrm{id}_X\) (the prescribed conditions of criterion (xi)(d)--(e), with \(o\) as the order record and \((\mathrm{id}_X)\) as \(L\)), a witness-schema annotation naming a
candidate outside \(P_1,\dots,P_6\) with \(\psi(C):\equiv r\in C\wedge o\notin
C\), where \(r\in C\) abbreviates the finite disjunction \(\bigvee_{d\in C}r=d\)
over the listed identifiers of \(C\) and \(o\notin C\) the conjunction
\(\bigwedge_{d\in C}o\neq d\), so that each instance \(\psi(C')\) is a closed
check-language formula, a check record carrying \(\psi(\{r\})\), \(\neg W_1(\{r\}),\dots,\neg
W_6(\{r\})\), \(\neg\psi(\{o\})\) and the finite conjunction over all
\(C'\subseteq\Raw{D_{\mathrm{scr}}}\) of \((W_1(C')\vee\dots\vee
W_6(C'))\rightarrow\neg\psi(C')\), a claim referencing \(r\) and that check
record whose statement is these conditions, and an honest audit certifying
it. The \(W_j\) atoms in these certificates are evaluated by the witness rows,
which read no certificate containing a \(W_j\) atom and, for \(W_6\), only
\(\mathrm{Audit}_0(D_{\mathrm{scr}})\), which excludes this audit; so the
certificates are well defined although the check record belongs to some
clusters \(C'\). Only \(o\) is law-bearing and no record supplies data for
\(W_1,W_2,W_3,W_5,W_6\), so \(W_i(C')\) holds only if \(o\in C'\), where
\(\psi(C')\) is false; no map, partial-map, transition or refinement record is
present, so the generated monoid is \(\{\mathrm{id}_X\}\). Recording \(c\) by
a map record would place it in the list of criterion (xi)(e), and \(r\) would
fail the screen. The assertion fields
of \(r\) are required for no witness, \(r\) carries no seed, begins no (R4)
chain and lies in the covered scope, so \(r\) satisfies (x) and (xi); the check record, which serves no witness and carries no seed, also receives (x) but not (xi), since it records no closure operation. Thus
(xi) is realizable in \FATCD{}, and its emptiness in
Theorem~\ref{thm:upper-bound-classification} comes from the hypothesis.

\emph{Open question.} Which assertion-bearing records admit a reduction in which the assertion is
itself a required field of an active cluster (label (i)), without
adjoining wrapper records such as the realization relation and one-class
packaging map above, or a trace reference, is left open and is
recorded with the endogenous-instrument meta-limit of
Sections~\ref{subsec:meta-limit} and~\ref{subsec:future-meta-limit}. The three
threat closures of
Propositions~\ref{prop:probability-closure}--\ref{prop:agency-closure} label
features of their recognized forms without the hypothesis; the classification
theorem below assumes it in general.
\end{remark}

Here and in the entry \emph{Closed unconditionally} of the working vocabulary in Section~\ref{subsec:objects-at-a-glance}, the phrase means that the stated recognized forms of Propositions~\ref{prop:probability-closure}--\ref{prop:agency-closure} receive labels without assuming the recognition hypothesis. It does not assert that hypothesis for every feature in those three threat classes; Theorems~\ref{thm:upper-bound-classification} and~\ref{thm:scoped-exact-six} retain it for general residual exhaustion.

Three ingredients combine in the next theorem. The alignment rules of
Proposition~\ref{prop:alignment-rules} fix which features each role can
claim. The threat closures of
Propositions~\ref{prop:probability-closure}--\ref{prop:agency-closure}
classify the recognized forms of probability, causality and agency. The residual scheme of
Definition~\ref{def:residual-scheme} supplies the labels. Under the
recognition hypothesis, the theorem assigns every residual feature a
label other than an unresolved threat or a seventh role: its proof
matches each recognized disposition to one of labels (i)--(ix). The
three preceding propositions classify their stated recognized forms
without that hypothesis.

The proof sorts a recognized feature by its disposition, following the
label table after Definition~\ref{def:residual-scheme}; a feature of any
raw-grammar or annotation kind receives the label of the disposition its
content has. The alternatives (R1)--(R5) of the recognition hypothesis are
written to correspond to the label criteria, so the proof matches each
alternative to a criterion and uses no substantive property of the particular
witness rows (Remark~\ref{rem:classification-scope}); which features the
hypothesis admits and which it excludes is described in
Remark~\ref{rem:recognition-open}.
Label~(viii) and the channel status \(\outsidescope\) of
Definition~\ref{def:channel-status} both concern content outside the
covered scope: the label sorts a single feature, and the status marks a
channel whose role content lies there.

\begin{theorem}[Upper-bound classification]\label{thm:upper-bound-classification}
Let \(D \in \FATCD{}\) satisfy the recognition hypothesis
(Definition~\ref{def:recognition-hypothesis}), let \(\Dec{D}\) be a
well-formed decomposition record of \(D\)
(Definition~\ref{def:decomposition-record}), and let \(r\) be a kind residual
of \(D\) (Definition~\ref{def:residual-scheme}: any content record or structural
annotation, including members of active clusters), a bookkeeping annotation of \(D\) (whether or not \(\Dec{D}\) assigns it to a channel), or a record residual, an
element of \(\Residual{D}\). Then some criterion among (i)--(ix) holds for \(r\) relative to \(\Dec{D}\); hence neither (x) nor (xi) is licensed for \(r\), and
every label that the label criteria after Definition~\ref{def:residual-scheme}
license for \(r\) relative to \(\Dec{D}\) is one of
\begin{enumerate}[label=(\roman*),topsep=2pt,itemsep=1pt]
  \item a declared constituent of the cluster that \(\Dec{D}\) names for an
    active \Pone{}, \Ptwo{}, \Pthree{}, \Pfour{}, \Pfive{}, or \Psix{}
    projection, in the sense of criterion (i) of Definition~\ref{def:residual-scheme};
  \item a shared constituent of two such projections;
  \item a presentation variant of an already classified feature;
  \item a level shift of an already classified feature;
  \item unabsorbed role-aligned content, covered-seed content with no
    certified witness \(W_i\), content checked in an honestly failed offer, or an activation-threshold refinement;
  \item an instrument-visibility artifact;
  \item an empirical-bridge annotation;
  \item an outside-domain structure;
  \item already-covered bookkeeping or annotation data, or role-neutral data
    without an assertion of its own.
\end{enumerate}
In particular, categories (x) and (xi) of
Definition~\ref{def:residual-scheme}, namely the unresolved-threat category and the recorded seventh-role
screen, are empty: no candidate
seventh-role feature in the covered scope remains as an unresolved threat, and
no feature in the covered scope passes the recorded seventh-role screen of
criterion (xi); here ``no unresolved threat'' means that some criterion among
(i)--(ix) holds for \(r\). Since seed clause (2) of
Definition~\ref{def:channel-status} accepts any one-variable macro-specification with a realization relation
whose recorded evaluations take both truth values, this certifies the form in
which \(r\) is recorded, not that its assertion concerns the role whose seed it
carries (Remark~\ref{rem:classification-scope}). \emph{Record residuals.} If \(r\) is an element of \(\Residual{D}\), every label licensed for it
lies in (iii)--(ix). \emph{Pointwise form.} The conclusion also holds pointwise, without the
recognition hypothesis for the rest of \(D\): for every \(D\in\FATCD{}\), every
well-formed \(\Dec{D}\) and every \(r\) as above that satisfies one of
(R1)--(R5), every label licensed for \(r\) lies in (i)--(ix), and in
(iii)--(ix) if \(r\in\Residual{D}\). For each fixed well-formed decomposition
record, alternatives (R1)--(R5) supply, respectively, a criterion among
(i)/(v), (v)/(vi)/(vii)/(ix-a)/(ix-b), (ix-c), (iii)/(iv), and (viii), as the
proof below establishes.
\end{theorem}

\begin{proof}
Alternatives (R1)--(R5) of Definition~\ref{def:recognition-hypothesis} were
written as record-independent forms of label criteria (i)--(ix) (paragraph
before Definition~\ref{def:recognition-hypothesis}); the proof checks, for a
fixed well-formed record, that each alternative meets one of those criteria.
Let \(D\in\FATCD{}\), fix a well-formed decomposition record
\(\Dec{D}\), and let \(r\) be a kind residual of \(D\)
(Definition~\ref{def:residual-scheme}), a bookkeeping annotation of \(D\), or an element of
\(\Residual{D}\) for that record. Fix any label
\(L\) that the label criteria after Definition~\ref{def:residual-scheme}
license for \(r\) relative to \(\Dec{D}\). The label \(L\) is drawn from
(i)--(xi); we show that, under the recognition hypothesis, it lies in
(i)--(ix), so that (x) and (xi) are empty.

By the label table, \(L\) is (x) only if no criterion (i)--(ix) holds for
\(r\), and \(L\) is (xi) only if \(r\) has the property in (x).

\emph{Conclusion.} Each recognition alternative of
Definition~\ref{def:recognition-hypothesis} supplies a label criterion among
(i)--(ix):
\begin{center}
\small
\begin{tabularx}{\linewidth}{@{}l>{\raggedright\arraybackslash}X@{}}
\toprule
Alternative & Label criterion supplied \\
\midrule
(R1) & (i) if criterion (i) holds; otherwise (v-a) for the first branch of (R1) and (v-c) for its failed-offer branch \\
(R2) & (vi) for visibility; (vii) for an admissible bridge claim (Definition~\ref{prop:honest-bookkeeping}) whose associated audit carries the five items of Definition~\ref{prop:empirical-bridge}; \\
     & (v-b) for a threshold annotation targeting a cluster named for an active projection and carrying only a threshold predicate besides typing, target and label fields, which reads at least one record, every record it reads being listed in its target field and satisfying (R1); \\
     & (ix-b) for remaining structural annotations satisfying (R2); (ix-a) for bookkeeping annotations satisfying the second clause of (R2) \\
(R3) & (ix-c) \\
(R4) & (iii) or (iv), the criterion of the terminal record being supplied by the rows for (R1), (R2), (R3) and (R5); no recursion occurs, since the terminal record satisfies (R1), (R2), (R3) or (R5) \\
(R5) & (viii) \\
\bottomrule
\end{tabularx}
\end{center}
Under the recognition hypothesis
(Definition~\ref{def:recognition-hypothesis}), every kind residual of \(D\), every bookkeeping annotation of \(D\) and every element of \(\Residual{D}\) for every well-formed decomposition record of \(D\), in particular \(r\), satisfies one of (R1)--(R5), so by the table a criterion among (i)--(ix)
holds for it; for (R4) the terminal record of the chain supplies it, and
chains are finite. This uses recognition of \(r\) alone and, for (R4), the
criterion that (R4) requires of the chain's terminal record, so it holds
pointwise in any \(D\in\FATCD{}\). Thus the criterion of (x), which requires that no
criterion (i)--(ix) hold, fails; since (xi) requires (x), it fails as well.
Hence \(L\) lies in (i)--(ix), and categories (x), the
unresolved-threat category, and (xi), the recorded seventh-role screen, are
empty. If \(r\in\Residual{D}\) for a decomposition record, \(L\) is neither (i)
nor (ii): by Definition~\ref{def:decomposition-record}, that set omits the
features already accounted for by an active projection, and (ii) requires
\(r\in C_i\cap C_j\subseteq A(\Dec{D})\); so every label licensed for a record
residual, relative to that record, lies in (iii)--(ix). The
emptiness of (x) and (xi) is the content of the recognition hypothesis, not a
consequence of the \FATCD{} definitions: by Remark~\ref{rem:recognition-open}
a finite audited conservation law or an emergent composite obstruction, if
unrecognized, would fall under (x), and under (xi) only if it also passes the recorded
seventh-role screen; Remark~\ref{rem:recognition-open} separately
exhibits an unrecognized normalization predicate in \FATCD{}.
\end{proof}

\begin{remark}[Scope of the theorem]\label{rem:classification-scope}
Theorem~\ref{thm:upper-bound-classification} is conditional on the recognition
hypothesis (Definition~\ref{def:recognition-hypothesis}) and scoped to \FATCD{}.
Alternative (R1) yields criterion (i) or (v); (R2) yields (v), (vi),
(vii), (ix-a) or (ix-b); (R3) yields (ix-c); (R4) yields (iii) or (iv); and
(R5) yields (viii). The alternatives are stated without reference to a
decomposition record; the theorem records that these
record-independent forms exclude (x). It does not derive that no finite audited feature requires a seventh role; its general conclusion assumes the recognition hypothesis. Propositions~\ref{prop:probability-closure}--\ref{prop:agency-closure} classify only their stated recognized forms without that hypothesis. It does
not claim that no seventh role can arise outside \FATCD{}, that the alignment
rules are exhaustive for richer domains, or that the recognition hypothesis holds
in general. Nor does it prove exact-six on its own: the decomposition/exhaustion
theorem of Section~\ref{sec:decomposition-exhaustion} remains necessary. Which further
classes of assertion-bearing records are recognized, and broader questions, are deferred to
Sections~\ref{subsec:meta-limit} and~\ref{sec:future-work}.
For each fixed well-formed \(\Dec{D}\), the proof matches alternatives
(R1)--(R5) to criteria (i)--(ix). This matching uses no substantive
property of the particular witness rows or seed clauses. It remains
valid for any finite family of check-language witness predicates and
seed clauses, provided the recognition alternatives, label criteria,
named seed fields, alignment and required-field conventions, and
active-cluster assignments are changed consistently. The number six enters
Theorem~\ref{thm:scoped-exact-six} through
Definition~\ref{def:decomposition-record}, and the particular predicates
through Theorem~\ref{thm:six-role-lower-bounds} and
Propositions~\ref{prop:probability-closure}--\ref{prop:agency-closure}. The (R1) and (R4) cases use that record's
active-cluster assignments and checked chains. The theorem proves that recognition excludes (x), using the table in its
proof. The converse does not follow from the label criteria: bookkeeping fields
may receive (ix) while substantive assertions require separate recognition,
and criteria (iii) and (iv) accept a chain whose terminal record lies in \(A(\Dec{D})\cup S(\Dec{D})\), for instance a member of an active or failed-offer cluster with an assertion field serving no witness check, which need not satisfy (R1)--(R5). Criteria (vi) and (ix-b)
impose the (R2) requirement that every record whose values a formula-valued prescribed field reads satisfy (R1), and each implies the structural-annotation branch of
(R2). The upper bound therefore depends on the stated
recognition hypothesis. Recognition is not necessary for the condition of Definition~\ref{def:scoped-exact-six-question}: adding to the nonclaim annotation of \(D_5\) a claim field with formula \(0+0=1\) gives \(D\in\FATCD{}\setminus\FATCD_{\mathrm{RH}}\) in which every well-formed record places that annotation (the only nonclaim annotation targeting \(C_5\), which the active \Pfive{} entry must assign) in \(A(\Dec{D})\), so no kind or record residual is licensed (x) or (xi). The proof uses recognition of kind residuals, bookkeeping annotations and record residuals. Since seed clause (2) accepts any one-variable
macro-specification with a realization relation whose recorded evaluations take both truth values, (R1)
certifies the form of a record, not that its assertion concerns the role
whose seed it carries.
\end{remark}
